\chapter{Calculabilité}\label{chapter-computability}

\section{Introduction à la calculabilité}
\subsection{Calculabilité, algorithmes et thèse de Church--Turing}
\subsection{Données finies, codages et fonctions partielles}

\section{Fonctions récursives et manipulation de programmes}
\subsection{Fonctions primitives récursives}
\subsection{Numérotation des fonctions primitives récursives}
\subsection{Auto-référence et théorème de récursion}
\subsection{Fonctions générales récursives}
\subsection{Le problème de l'arrêt}

\section{Machines de Turing}
\subsection{Définition et fonctionnement des machines de Turing}
\subsection{Calcul par machine de Turing et simulation}
\subsection{Machines de Turing et indécidabilité}

\section{Décidabilité, semi-décidabilité et réductions}
\subsection{Décidabilité et semi-décidabilité}
\subsection{Le théorème de Rice}
\subsection{Réductions et problèmes indécidables}

\section{Le \texorpdfstring{$\lambda$}{lambda}-calcul non typé}
\subsection{Syntaxe, variables et substitution}
\subsection{Réduction et stratégies d'évaluation}
\subsection{Entiers de Church et fonctions représentables}

\section{Conclusion}\label{section-computability-conclusion}


\chapter{Typage}\label{chapter-typing}

\section{Généralités sur le typage}
\subsection{Rôle et variétés des systèmes de typage}
\subsection{Constructions usuelles sur les types}
\subsection{Polymorphisme et tâches d'un système de typage}
\subsection{Autres usages et exemples}
\subsection{Typage, terminaison et correspondance de Curry--Howard}

\section{Le \texorpdfstring{$\lambda$}{lambda}-calcul simplement typé}
\subsection{Syntaxe et règles de typage}
\subsection{Propriétés élémentaires du typage}
\subsection{Substitution, réduction et normalisation}
\subsection{Reconstruction des termes et problèmes algorithmiques}

\section{Le calcul propositionnel intuitionniste}
\subsection{Formules et déduction naturelle}
\subsection{Théorèmes et non-théorèmes}
\subsection{Logique classique}
\subsection{Interprétation de Brouwer--Heyting--Kolmogorov}
\subsection{Correspondance de Curry--Howard}
\subsection{Calcul des séquents et élimination des coupures}

\section{Le call/cc et la logique classique}
\subsection{Continuations}
\subsection{Continuation Passing Style}

\section{Sémantique du calcul propositionnel intuitionniste}
\subsection{Sémantiques, correction et complétude}
\subsection{Cadres de Kripke}
\subsection{Sémantique topologique}
\subsection{Réalisabilité propositionnelle}
\subsection{Sémantique des problèmes finis}

\section{Inférence de type à la Hindley--Milner}
\subsection{Désannotation et problème de l'inférence}
\subsection{Collection et résolution des contraintes}
\subsection{Polymorphisme du \texttt{let} et restriction de valeur}


\chapter{Introduction aux quantificateurs}
\label{chapter-quantifiers}

\section{Les quantificateurs et les systèmes logiques}
\subsection{Les quantificateurs comme conjonctions et disjonctions en famille}
\subsection{Interprétation de Brouwer--Heyting--Kolmogorov}
\subsection{Introduction et élimination de $\forall$ et $\exists$}
\subsection{Preuves, termes d'individus et notations de Curry--Howard}
\subsection{Ce que l'on peut quantifier : systèmes de types et $\lambda$-cube}
\subsection{Quantification existentielle, types sommes et imprédicativité}

\section{Logique du premier ordre}
\subsection{Individus, prédicats et syntaxe des formules}
\subsection{Règles de déduction naturelle et conditions de fraîcheur}
\subsection{Exemples de démonstrations et notation des preuves}
\subsection{Le monde des individus et le monde logique}
\subsection{Égalité et schémas d'axiomes}
\subsection{Curry--Howard pour la logique du premier ordre}

\section{Arithmétique du premier ordre}
\subsection{Termes arithmétiques et axiomes de Peano}
\subsection{Arithmétique de Heyting et arithmétique de Peano}
\subsection{Exemples de preuves par récurrence}
\subsection{Propriétés arithmétiques et formalisation des calculs}
\subsection{Programmes extraits des preuves et propriétés d'existence}

\section{Formalisation des démonstrations et théorème de Gödel}
\subsection{Codage des formules, démonstrations et calculs}
\subsection{Démontrabilité, semi-décidabilité et problème de l'arrêt}
\subsection{L'argument de Gödel, Rosser et Turing}
\subsection{Incomplétude de l'arithmétique de Peano}
\subsection{Consistance, interprétation et portée du théorème}
